Network operators must infer path headroom from measurements that are often intrusive, incomplete, or too coarse to explain why performance has changed. Modern services run over networks whose traffic mix, topology, routing state, queue occupancy, and failure modes change continuously. A central version of this problem is remaining path capacity, meaning how much additional traffic a path can absorb before delay or user-visible performance starts to degrade. Overestimating this headroom risks congestion and service disruption. Underestimating it wastes capacity and can lead to unnecessary routing or provisioning decisions.

Delay is one of the most accessible signals for this question, but it is often compressed too early. Operational systems commonly report mean round-trip time, high-percentile latency, link utilisation, or threshold alarms. These summaries are useful, but each is a scalar view of a richer object. A path whose mean delay remains stable may still be developing a heavier tail. A second delay mode may appear when traffic is rerouted through a longer path or when queues begin to form intermittently. Two path states can have similar percentile values while having different distributional shapes and therefore different operational implications. Treating delay as a distribution rather than as a single number keeps this information available.

The central question is whether receiver-visible delay samples contain enough information to detect when a path leaves its low-load delay regime. The intended monitor observes packet delays, reconstructs the empirical delay distribution, decides when that reconstruction is stable, and then uses distributional change to infer whether increasing offered load has pushed the path into a different regime. The configured capacity of a simulated link may be used after an experiment to assess error, but it is not available to the monitoring algorithm. The method is designed to reason from observations, not from knowledge of the network configuration.

Existing monitoring systems tend to fall into three groups: active tools that infer available bandwidth by injecting load, telemetry systems that expose queue or hop state from inside the network, and threshold systems that detect anomalies through scalar summaries. The gap addressed here is narrower. The question is whether receiver-visible delay distributions alone can identify the point at which a path leaves its low-load operating regime, and whether that decision can be made without ground-truth capacity or queue state.

\section{Aims and contributions}

The aim of the project is to test whether delay distributions can be treated as a practical monitoring object for estimating a capacity-relevant boundary. This is narrower than general anomaly detection and weaker than direct queue telemetry, but it is chosen because it only assumes receiver-visible delay samples. A successful result would show that distributional monitoring can provide useful evidence before a coarse endpoint alarm is triggered, without requiring the monitor to know the underlying delay family or the true link capacity.

The work is organised around four sub-goals. The first is to establish the statistical basis for learning an unknown delay distribution from samples. The second is to decide when such a reconstruction is stable enough to be used online, where ground truth is unavailable. The third is to turn the stopped reconstruction into an identity decision between operating states. The fourth is to test whether that decision remains meaningful as the network setting moves from direct sampling to a single link, longer paths, and separable multi-path mixtures.

These goals are expressed as three research questions in Table \ref{tab:research-questions}. They are intentionally sequential. A capacity-boundary decision made from delay distributions is only meaningful if the distribution can first be reconstructed, if the reconstruction can be stopped without access to truth, and if the resulting comparison identifies more than ordinary low-load drift.

\begin{table}[htbp]
\centering
\small
\begin{tabular}{p{0.12\textwidth}p{0.78\textwidth}}
\toprule
Question & Focus \\
\midrule
RQ1 & Can an unknown delay distribution be reconstructed from receiver-side samples within a total variation distance target of \(\delta=0.05\)? \\
RQ2 & Can stability between consecutive reconstructions act as an online stopping signal when the true distribution is unavailable? \\
RQ3 & Can changes in stopped delay distributions identify a capacity-relevant boundary as offered load increases, including when the path becomes longer or the receiver sees a mixture of path behaviours? \\
\bottomrule
\end{tabular}
\caption{Research questions addressed by the evaluation.}
\label{tab:research-questions}
\end{table}

The chapter order follows this development. The use case first defines the monitoring pipeline and the intended capacity decision. The statistical chapter then tests whether an unknown delay distribution can be reconstructed from samples. Adaptive stopping tests whether that reconstruction can be trusted without seeing ground truth. The network chapters then use the same stopped distributions to refine the operational decision rule across single-link, multi-hop, and multi-path settings.

The claim is bounded. The report does not claim to measure exact physical capacity. The output is a distributional boundary: the first offered-load region where the observed delay distribution begins moving away from the low-load state much faster than its previous trend. Configured link capacity is used only for post-hoc evaluation in simulation.

The contributions are summarised in Table \ref{tab:intro-contributions}. The main research result is the replacement of first distributional change with a response-rate criterion. Capacity pressure is better identified by the point where total variation distance begins increasing much faster than its previous low-load trend, rather than by the first point where two empirical distributions differ.

\begin{table}[htbp]
\centering
\small
\begin{tabular}{p{0.12\textwidth}p{0.78\textwidth}}
\toprule
Contribution & Summary \\
\midrule
C1 & An empirical calibration of non-parametric delay-distribution reconstruction against a finite-domain total variation distance learning reference. \\
C2 & An adaptive stopping evaluation that shows where consecutive-estimate stability is reliable, and where false stability causes premature stopping. \\
C3 & A capacity-boundary evaluation showing that first distributional change is too sensitive, while a response-rate criterion gives a repeatable boundary under constant-rate single-link and multi-hop conditions and remains interpretable under bursty and heterogeneous stress tests. \\
C4 & A constructed multi-path modal analysis showing that separable receiver-side mixture distributions should be decomposed before capacity association, because whole-mixture TVD can dilute or misplace the stressed path response. \\
\bottomrule
\end{tabular}
\caption{Main contributions of the work.}
\label{tab:intro-contributions}
\end{table}

\section{Network performance context}

The difficulty of this task begins with the structure of the networks being monitored. Large data-centre networks are designed to provide high bisection bandwidth, path diversity, and efficient resource pooling. Fat-tree and Clos-derived designs made it possible to construct large fabrics from commodity switching elements, while systems such as VL2 and Jupiter demonstrate how multi-stage topologies and centralised control can support large-scale cloud services \cite{alFaresFatTree,greenbergVl2,singhJupiter}. These designs improve aggregate capacity, but they also create a measurement problem. Traffic can traverse multiple equal-cost paths, failures can be masked by rerouting, and a performance problem may be visible only as a change in latency, queueing, or flow completion time.

The transport layer adds another reason to care about delay distributions. Data Center TCP (DCTCP), a variant of Transmission Control Protocol, was motivated by the coexistence of latency-sensitive short flows and throughput-hungry background flows, where queue build-up from large flows can harm foreground latency even before severe congestion dominates \cite{alizadehDctcp}. This setting shows why a scalar utilisation counter can be misleading. A link may not appear fully saturated on average, while packet delay begins to show a longer tail or a second operating mode. Remaining capacity is therefore not only a question of nominal link rate. It is a question about when the path's observed behaviour leaves the regime in which it can support additional traffic without unacceptable delay.

Traditional network-management mechanisms provide the historical baseline. Simple Network Management Protocol (SNMP) defines a manager-agent architecture for reading and controlling management information on network devices \cite{harringtonSnmp}. Flow-level systems such as NetFlow export records describing unidirectional Internet Protocol flows, while sampled flow (sFlow) uses packet and counter sampling to support continuous traffic monitoring in high-speed switched and routed networks \cite{claiseNetflow,phaalSflow}. These mechanisms are valuable because they scale to large deployments and have become part of operational practice. Their limitation for this project is their abstraction. Interface counters and flow records are well suited to accounting, capacity planning, and coarse anomaly detection, but they do not directly provide a stable empirical distribution of packet delay on a path.

The key design tension is therefore granularity against overhead. Fine-grained monitoring reveals short-lived queueing and path behaviour that coarse counters miss. However, observing every packet, carrying per-hop metadata, or injecting synthetic probe traffic can consume bandwidth, switch resources, or operational effort. The approach studied here occupies the low-assumption side of that trade-off. It treats delay samples visible at the measurement edge as the primary evidence, then asks how much can be inferred statistically before requiring richer network instrumentation.

\section{Active monitoring}

Active measurement estimates path properties by sending measurement traffic. This family is important because the capacity use case is naturally tempting to frame as active probing. Pathload estimates available bandwidth by sending self-loading periodic streams and detecting increasing one-way delay trends when the probing rate exceeds the available bandwidth \cite{jainDovrolisPathload}. PathChirp improves efficiency by sending exponentially spaced probe trains, extracting information from the rate range covered by each chirp rather than repeatedly probing at a single rate \cite{ribeiroPathchirp}. Packet-pair and packet-train dispersion methods infer capacity or available bandwidth from the spacing of probe packets after they traverse a bottleneck, while Spruce was proposed as a lightweight alternative evaluated against Pathload and Initial Gap Increasing or Packet Transmission Rate methods over wide-area paths \cite{dovrolisPacketDispersion,straussSpruce}.

These methods provide a clear conceptual foundation because delay trends and packet spacing contain information about spare capacity. They are also attractive because they can be deployed from end hosts without requiring internal switch support. The difficulty is that the measurement traffic is part of the experiment. A probe stream must be large enough to reveal congestion, but that same stream can perturb the path being measured. The problem is more acute when the operational question is how much headroom remains. Near the boundary, extra traffic is not merely a measurement cost; it is the mechanism that may push the path into a worse regime.

Coverage-oriented active systems make a different trade-off. Pingmesh continuously measures latency between servers at large scale, while RD-Probe combines randomised and deterministic probing to cover physical components in complex data-centre networks \cite{guoPingmesh,dingRdProbe}. These systems are valuable because they show how active measurement can make hidden network states visible across a large deployment. Their objective, however, is coverage and troubleshooting rather than learning a full delay distribution at a capacity boundary.

Active tools also usually compress their output into a rate estimate or a reachability result. That is appropriate when the goal is to report available bandwidth or path coverage, but it discards some of the signal that appears before, during, and after the rate boundary. A path approaching congestion may first show changes in skew, tail mass, or multimodality rather than a single monotonic delay trend. The capacity use case in this report still changes offered load experimentally, because the goal is to validate whether a capacity boundary can be detected. The monitoring decision, however, is made from the empirical delay distributions observed at each operating point rather than from a direct probe-rate equation. The eventual measurement object is the distributional state of the path under the current load, rather than the probe stream itself.

\section{Passive monitoring}

Passive monitoring avoids injecting dedicated measurement traffic and instead uses traffic already visible to the monitor. For the remaining-capacity use case, this is the more natural target interface. If a monitor can infer a capacity-relevant change from observed delays, then it can run continuously without consuming the headroom it is trying to estimate. Passive monitoring also aligns with operational settings where active probes are undesirable because of policy restrictions, customer traffic sensitivity, or the risk of perturbing short-lived congestion events.

Large-scale passive and hybrid systems show the breadth of this design space. Fathom passively samples remote procedure calls in Google's production fleet and decomposes application latency into client, network, and server components \cite{qureshiFathom}. It is powerful because it links network performance to application-level work and aggregates sampled measurements into distributions. Its measurement interface is nevertheless service-specific and relies on host and stack instrumentation. PIRATE addresses a different visibility constraint by estimating response latency when only one traffic direction is visible under asymmetric routing, using causal request pairs and packet timing gaps \cite{shobhanaPirate}. This demonstrates that packet timing can reveal latency information even when visibility is incomplete. The inferred quantity, however, is response latency from causal structure rather than a path delay distribution used to estimate capacity headroom.

Other systems use passive packet capture or packet histories to improve diagnosis. NetSight records packet histories so that operators can query how packets moved through a network, while Everflow traces selected packets in large data-centre networks using match-and-mirror functionality and analysis servers \cite{handigolNetSight,zhuEverflow}. These systems offer packet-level diagnostic power, but the price is a heavier measurement substrate. They need packet capture, mirroring, collection infrastructure, or packet histories. For debugging and root-cause analysis that cost can be justified. It is still a different operating point from a lightweight edge monitor that estimates whether the delay distribution of a path has changed.

Fine-grained passive systems also show why sub-second and microsecond effects cannot be ignored. Trumpet uses end-host triggers to detect network-wide events at millisecond timescales while inspecting packets at line rate \cite{moshrefTrumpet}. \(\mu\)Mon compresses microsecond-level flow-rate curves using wavelets and replays congestion events to recover short-lived dynamics \cite{zhengMuMon}. These systems reinforce the motivation for moving beyond slow counters. They also highlight a distinction in scope. Their primary objects are events, triggers, or rate curves. The object here is a learned delay distribution, and the central question is whether its stable shape changes as load approaches a capacity limit.

\section{Programmable and in-band telemetry}

Programmable data planes and in-band network telemetry (INT) provide the strongest visibility when the infrastructure supports them. The P4 INT specification allows packets to carry switch-local metadata such as hop latency, queue occupancy, timestamps, and egress utilisation as they traverse the network \cite{p4IntSpec}. Marple proposes a performance query language and switch hardware primitive for line-rate aggregations such as moving averages over queueing latency per flow \cite{narayanaMarple}. ConQuest uses programmable-switch data structures to identify flows that contribute significantly to queue build-up, reporting high precision at millisecond timescales with a small memory footprint \cite{chenConQuest}. UnivMon provides a more general sketch-based framework for estimating network-wide monitoring metrics such as heavy hitters and entropy using programmable measurement \cite{liuUnivMon}.

These systems are compelling because they observe the network close to where queueing occurs. If a switch can report queue occupancy, hop latency, and per-flow contributions, then capacity-related diagnosis can be much more direct. The disadvantage is that this assumes cooperation from the network core. The monitor needs programmable devices, switch resources, telemetry headers, mirrored traffic, or collection infrastructure. That assumption is not always available. It is particularly restrictive for a method intended to reason from path observations rather than from device-local state.

In-band network telemetry also has an overhead problem. Carrying hop-by-hop metadata in packets increases packet size and can affect application performance. Probabilistic in-band network telemetry (PINT) responds to this by encoding telemetry information across packets so that per-packet overhead can be bounded, even as low as one bit for some queries \cite{benBasatPint}. The existence of PINT is itself evidence that raw INT is not free. Reducing overhead is possible, but it introduces a telemetry design problem involving what information should be encoded, how accurately it should be recovered, and which packets should carry it.

The approach in this report is weaker in instrumentation but stronger in deployability. It does not claim to replace INT when switch-local queue state is available. Instead, it targets the complementary case where such support is absent or too expensive. The comparison is therefore not that passive distributional monitoring reveals more than programmable telemetry in an instrumented fabric. Rather, it asks whether a useful capacity signal remains when the available evidence is only the end-to-end delay samples that the monitor can observe.

\section{Threshold-based monitoring}

Operational monitoring often turns continuous measurements into alarms through thresholds. Static thresholds are easy to understand and explain. A link utilisation above a configured percentage, a latency percentile above a service-level objective, or an error rate above a fixed tolerance gives an operator a direct signal. This clarity is one reason threshold-based systems remain common. The weakness is that the threshold is usually imposed on a scalar metric. A single number must stand in for the behaviour of a complex process.

More sophisticated threshold systems adapt to time-series structure. Brutlag integrated Holt-Winters forecasting into network service monitoring so that expected seasonal behaviour could be modelled and deviations could trigger alarms \cite{brutlagAberrant}. Barford, Kline, Plonka, and Ron used signal analysis and wavelets to expose network traffic anomalies across multiple timescales \cite{barfordSignalAnalysis}. Lakhina, Crovella, and Diot used principal component analysis to separate normal and anomalous subspaces in network-wide traffic matrices, diagnosing volume anomalies with low false alarm rates \cite{lakhinaDiagnosing}. These systems all improve on naive fixed thresholds by learning or decomposing normal traffic structure.

The limitation remains that an alarm is usually driven by deviation in an engineered summary. In service-level monitoring, Sommers, Barford, Duffield, and Ron study how to monitor compliance guarantees using distribution-free techniques for quantiles \cite{sommersSla}. That line of work is closer to the problem here because delay guarantees are statistical rather than purely average-based. However, quantile compliance is still narrower than comparing full empirical distributions. Two delay distributions can agree at the monitored quantile while differing elsewhere in the body or tail. Conversely, a small but persistent movement of probability mass can be operationally meaningful before a chosen threshold is crossed.

The proposed monitor keeps the interpretability of a threshold, but moves the threshold onto a distributional distance. Instead of asking whether a single latency value is too high, the monitor asks whether the current stopped empirical distribution has moved too far from an empirical operating reference. The threshold is still necessary, because operational decisions require a boundary. The difference is that the boundary is applied after preserving shape information rather than before it.

\section{Distributional monitoring}

Distributional ideas already appear in several areas of network measurement. Lakhina, Crovella, and Diot showed that traffic feature distributions can reveal and classify anomalies that are difficult to detect from volume alone \cite{lakhinaDistributions}. Nychis, Sekar, Andersen, Kim, and Zhang later evaluated entropy-based anomaly detection across multiple traffic distributions, showing both the promise of distribution-derived metrics and the care needed when selecting features and interpreting entropy changes \cite{nychisEntropy}. These works are important because they demonstrate that the distribution of traffic features can carry diagnostic information hidden from scalar volume measurements.

Delay-specific distributional work makes the connection even clearer. Network delay tomography estimates internal link delay distributions from end-to-end measurements, using statistical inference to recover hidden delay components without direct internal instrumentation \cite{tsangDelayTomography}. Packet Doppler monitors packet shifts and argues that passive distributional monitoring can detect burst events missed by active probing \cite{qiuPacketDoppler}. Empirical delay-modelling work has also used families such as lognormal and Weibull to capture skewed Internet delay and traffic behaviour \cite{karakasDelay,downeyLognormal,arfeenWeibull}. These papers justify treating delay shape as meaningful rather than as noise around a mean.

The remaining gap is not that distributions have been ignored. The gap is that distributional monitoring has usually been used for anomaly detection, tomography, or model fitting, rather than for estimating remaining path capacity through repeated comparison of stable delay distributions. Anomaly detection asks whether something unusual happened. Tomography asks where hidden delay components are located. Parametric modelling asks which family can describe observed measurements. The present work asks a narrower operational question about when a path leaves the delay distribution associated with its low-load operating state as offered load increases.

The statistical design follows from that distinction. If a monitor compares unstable empirical distributions, it risks interpreting sampling noise as congestion. If it waits for a worst-case sample-complexity bound every time, it may be too slow for online use. The project therefore combines distribution reconstruction, adaptive stopping, and distribution identity comparison. Reconstruction asks whether the empirical distribution is accurate enough in controlled settings. Adaptive stopping asks when enough samples have been observed without access to ground truth. Identity comparison then asks whether two stopped distributions are sufficiently different to indicate a capacity-relevant change.

\section{Project position}

The contribution is a statistical and experimental argument for treating delay distributions as first-class path characteristics. It is narrower than general network anomaly detection because the motivating question is remaining capacity. It is weaker than programmable telemetry because it does not assume switch-local queue state, in-band metadata, or packet histories. It is less intrusive than active available-bandwidth probing because the monitor's decision object is the observed delay distribution, not a probe stream that must itself load the path. These restrictions are deliberate. They define a measurement interface that is realistic for passive monitoring and still rich enough to test whether capacity pressure leaves a detectable distributional signature.

The work is organised around three statistical tasks. The first task is reconstruction. Given samples from an unknown delay process, the monitor estimates a probability mass function or a smoothed density without being told whether the data came from a normal, lognormal, Weibull, binomial, Zipfian, or irregular distribution. The second task is stopping. A monitor cannot wait for a worst-case theoretical number of samples every time, so an online rule is needed to decide when the current distribution estimate has stabilised. The third task is identification. Once a stable distribution has been reconstructed, the monitor compares it with another stable distribution to decide whether the path has changed enough to matter operationally.

The experimental scope is intentionally controlled. Direct sampling first isolates the statistical reconstruction problem from network effects. The next stage applies the same reconstruction and stopping logic to a simple Network Simulator 3 (NS-3) topology with controlled offered load. Later stages then increase topology complexity through multi-hop paths and multi-path mixtures. This staged design makes the evaluation diagnostic. If a network experiment fails, the likely source of failure can be separated from basic distribution learning and traced instead to the stopping rule, the network mechanism, or the distributional decision boundary.

The central claim is therefore not that delay distributions solve all monitoring problems. They cannot identify the responsible flow as directly as ConQuest, expose hop metadata as directly as INT, or produce a traditional available-bandwidth estimate in the same form as Pathload. Their value is different. A delay distribution can be observed from the edge, compared with statistical guarantees, and interpreted as a behavioural fingerprint of a path under load. If this fingerprint changes consistently near capacity, then passive distributional monitoring gives operators a practical way to reason about headroom without assuming full control over the network interior.
